\documentclass{ximera}
%verbeterd 9/9/26
\title{Vraagstukken elektrische krachtwerking: goniometrische figuren}
\author{Daan Lipkens}

\begin{document}
\begin{abstract}
    Oefeningen op coulombkracht met goniometrische figuren (enkel pythagoras), van makkelijk naar moeilijk.
\end{abstract}
\maketitle

%———————————————————————————————————————
% OEFENING 1 (makkelijkst): rechthoekige driehoek, loodrechte krachten (Pythagoras)
\begin{exercise}
    Drie puntladingen bevinden zich in een assenstelsel. Lading $Q_{1} = -5{,}0 \cdot 10^{-6} \, \mathrm{C}$ ligt in de oorsprong $(0{,}00 \, \mathrm{m} ; 0{,}00 \, \mathrm{m})$. Lading $Q_{2} = +4{,}0 \cdot 10^{-6} \, \mathrm{C}$ ligt op de $x$-as op een afstand van $0{,}40 \, \mathrm{m}$ van de oorsprong. Lading $Q_{3} = +3{,}0 \cdot 10^{-6} \, \mathrm{C}$ ligt op de $y$-as op een afstand van $0{,}30 \, \mathrm{m}$ van de oorsprong.

    Bereken de grootte van de resulterende netto kracht op $Q_{1}$.

    \begin{center}
        \begin{tikzpicture}[scale=1]
            \useasboundingbox (-1.5, -1.5) rectangle (5.5, 4.5);
            \coordinate (A) at (0,0);
            \coordinate (B) at (4,0);
            \coordinate (C) at (0,3);

            % Driehoek en assen tekenen
            \draw[dashed, gray] (A) -- (B) node[midway,below] {$0{,}40 \, \mathrm{m}$};
            \draw[dashed, gray] (A) -- (C) node[midway,left] {$0{,}30 \, \mathrm{m}$};
            \draw[dashed, lightgray] (C) -- (B);

            % Ladingen labelen
            \tkzLabelPoint[below left](A){$Q_{1}$}
            \tkzLabelPoint[below right](B){$Q_{2}$}
            \tkzLabelPoint[above left](C){$Q_{3}$}

            % Ladingen tekenen
            \draw[charge-] (A) circle (\R) node[scale=1] {$-$};
            \draw[charge+] (B) circle (\R) node[scale=1] {$+$};
            \draw[charge+] (C) circle (\R) node[scale=1] {$+$};
        \end{tikzpicture}
    \end{center}

    \[ F_{\text{res},1} = \answer[tolerance=0.05,onlineshowanswerbutton]{1.9} \, \mathrm{N} \]

    \begin{hint}
        Teken eerst de krachten $\vec{\mathbf{F}}_{21}$ en $\vec{\mathbf{F}}_{31}$. Omdat $Q_1$ negatief is en de andere twee positief zijn, zijn het beide aantrekkingskrachten.
    \end{hint}
    \begin{hint}
        Bereken de groottes van $F_{21}$ en $F_{31}$ met de wet van Coulomb. Omdat de ladingen op de $x$- en $y$-as liggen, staan de vectoren perfect loodrecht ($90^\circ$) op elkaar.
    \end{hint}
    \begin{hint}
        Als de hoek tussen twee vectoren $90^\circ$ is, vereenvoudigt de cosinusregel zich tot de stelling van Pythagoras: $F_{\text{res}} = \sqrt{F_{21}^2 + F_{31}^2}$.
    \end{hint}

    \begin{solution}\nl
        \underline{Gegevens:}
        \begin{align*}
            & Q_{1} = -5{,}0 \cdot 10^{-6} \, \mathrm{C} & & \\
            & Q_{2} = +4{,}0 \cdot 10^{-6} \, \mathrm{C} & & r_{12} = 0{,}40 \, \mathrm{m}\\
            & Q_{3} = +3{,}0 \cdot 10^{-6} \, \mathrm{C} & & r_{13} = 0{,}30 \, \mathrm{m}
        \end{align*}

        \underline{Gevraagd:} $F_{\text{res},1}$\\

        \underline{Oplossing:}\\
        We berekenen eerst de groottes van de twee krachten die inwerken op lading $Q_{1}$.
        \begin{align*}
            F_{21} &= k\frac{\lvert Q_{2}\rvert\cdot\lvert Q_{1}\rvert}{(r_{12})^{2}} \\
            &= \frac{8{,}99 \cdot 10^9 \, \mathrm{\frac{N \cdot m^2}{C^2}} \cdot 4{,}0 \cdot 10^{-6} \, \mathrm{C} \cdot 5{,}0 \cdot 10^{-6} \, \mathrm{C}}{(0{,}40 \, \mathrm{m})^{2}} \\
            &= 1{,}1238 \, \mathrm{N}
        \end{align*}
        Deze kracht is een aantrekking en wijst dus langs de $x$-as naar rechts.

        \begin{align*}
            F_{31} &= k\frac{\lvert Q_{3}\rvert\cdot\lvert Q_{1}\rvert}{(r_{13})^{2}} \\
            &= \frac{8{,}99 \cdot 10^9 \, \mathrm{\frac{N \cdot m^2}{C^2}} \cdot 3{,}0 \cdot 10^{-6} \, \mathrm{C} \cdot 5{,}0 \cdot 10^{-6} \, \mathrm{C}}{(0{,}30 \, \mathrm{m})^{2}} \\
            &= 1{,}4983 \, \mathrm{N}
        \end{align*}
        Deze kracht is eveneens een aantrekking en wijst dus langs de $y$-as naar boven.

        Omdat de vectoren op de assen liggen, is de hoek hiertussen exact $90^\circ$. We kunnen de grootte van de resulterende vector daarom eenvoudig met de stelling van Pythagoras berekenen:
        \begin{align*}
            F_{\text{res},1} &= \sqrt{F_{21}^{2} + F_{31}^{2}} \\
            &= \sqrt{(1{,}1238 \, \mathrm{N})^{2} + (1{,}4983 \, \mathrm{N})^{2}} \\
            &= \sqrt{1{,}2628 + 2{,}2450} \\
            &= \sqrt{3{,}5078} \\
            &= 1{,}8729 \, \mathrm{N} \simeq 1{,}9 \, \mathrm{N}
        \end{align*}
        De resulterende kracht is ongeveer $1{,}9 \, \mathrm{N}$ groot.
    \end{solution}
\end{exercise}
%———————————————————————————————————————


%———————————————————————————————————————
% OEFENING EXTRA 1: Rechthoekige driehoek (Pythagoras)
\begin{exercise}
    Drie puntladingen bevinden zich in een assenstelsel. Lading $Q_{1} = -5{,}0 \cdot 10^{-6} \, \mathrm{C}$ ligt in de oorsprong $(0{,}00 \, \mathrm{m} ; 0{,}00 \, \mathrm{m})$. Lading $Q_{2} = +4{,}0 \cdot 10^{-6} \, \mathrm{C}$ ligt op de $x$-as op een afstand van $0{,}40 \, \mathrm{m}$ van de oorsprong. Lading $Q_{3} = +3{,}0 \cdot 10^{-6} \, \mathrm{C}$ ligt op de $y$-as op een afstand van $0{,}30 \, \mathrm{m}$ van de oorsprong.

    Bereken de grootte van de resulterende netto kracht op $Q_{1}$.

    \begin{center}
        \begin{tikzpicture}[scale=1]
            \useasboundingbox (-1.5, -1.5) rectangle (5.5, 4.5);
            \coordinate (A) at (0,0);
            \coordinate (B) at (4,0);
            \coordinate (C) at (0,3);
            
            % Driehoek en assen tekenen
            \draw[dashed, gray] (A) -- (B) node[midway,below] {$0{,}40 \, \mathrm{m}$};
            \draw[dashed, gray] (A) -- (C) node[midway,left] {$0{,}30 \, \mathrm{m}$};
            \draw[dashed, lightgray] (C) -- (B);
        
            % Ladingen labelen
            \tkzLabelPoint[below left](A){$Q_{1}$}
            \tkzLabelPoint[below right](B){$Q_{2}$}
            \tkzLabelPoint[above left](C){$Q_{3}$}
        
            % Ladingen tekenen
            \draw[charge-] (A) circle (\R) node[scale=1] {$-$};       
            \draw[charge+] (B) circle (\R) node[scale=1] {$+$};
            \draw[charge+] (C) circle (\R) node[scale=1] {$+$};
        \end{tikzpicture}
    \end{center}
    
    \[ F_{\text{res},1} = \answer[tolerance=0.05,onlineshowanswerbutton]{1.87} \, \mathrm{N} \]

    \begin{hint}
        Teken eerst de krachten $\vec{\mathbf{F}}_{21}$ en $\vec{\mathbf{F}}_{31}$. Omdat $Q_1$ negatief is en de andere twee positief zijn, zijn het beide aantrekkingskrachten.
    \end{hint}
    \begin{hint}
        Bereken de groottes van $F_{21}$ en $F_{31}$ met de wet van Coulomb. Omdat de ladingen op de $x$- en $y$-as liggen, staan de vectoren perfect loodrecht ($90^\circ$) op elkaar.
    \end{hint}
    \begin{hint}
        Als de hoek tussen twee vectoren $90^\circ$ is, vereenvoudigt de cosinusregel zich tot de stelling van Pythagoras: $F_{\text{res}} = \sqrt{F_{21}^2 + F_{31}^2}$.
    \end{hint}

    \begin{solution}\nl
        \underline{Gegevens:}
        \begin{align*}
            & Q_{1} = -5{,}0 \cdot 10^{-6} \, \mathrm{C} & & \\
            & Q_{2} = +4{,}0 \cdot 10^{-6} \, \mathrm{C} & & r_{12} = 0{,}40 \, \mathrm{m}\\
            & Q_{3} = +3{,}0 \cdot 10^{-6} \, \mathrm{C} & & r_{13} = 0{,}30 \, \mathrm{m}
        \end{align*}

        \underline{Gevraagd:} $F_{\text{res},1}$\\ 

        \underline{Oplossing:}\\
        We berekenen eerst de groottes van de twee krachten die inwerken op lading $Q_{1}$.
        \begin{align*}
            F_{21} &= k\frac{\lvert Q_{2}\rvert\cdot\lvert Q_{1}\rvert}{(r_{12})^{2}} \\
            &= \frac{8{,}99 \cdot 10^9 \, \mathrm{\frac{N \cdot m^2}{C^2}} \cdot 4{,}0 \cdot 10^{-6} \, \mathrm{C} \cdot 5{,}0 \cdot 10^{-6} \, \mathrm{C}}{(0{,}40 \, \mathrm{m})^{2}} \\
            &= 1{,}1238 \, \mathrm{N}
        \end{align*}
        Deze kracht is een aantrekking en wijst dus langs de $x$-as naar rechts.
        
        \begin{align*}
            F_{31} &= k\frac{\lvert Q_{3}\rvert\cdot\lvert Q_{1}\rvert}{(r_{13})^{2}} \\
            &= \frac{8{,}99 \cdot 10^9 \, \mathrm{\frac{N \cdot m^2}{C^2}} \cdot 3{,}0 \cdot 10^{-6} \, \mathrm{C} \cdot 5{,}0 \cdot 10^{-6} \, \mathrm{C}}{(0{,}30 \, \mathrm{m})^{2}} \\
            &= 1{,}4983 \, \mathrm{N}
        \end{align*}
        Deze kracht is eveneens een aantrekking en wijst dus langs de $y$-as naar boven.

        Omdat de vectoren op de assen liggen, is de hoek hiertussen exact $90^\circ$. We kunnen de grootte van de resulterende vector daarom eenvoudig met de stelling van Pythagoras berekenen:
        \begin{align*}
            F_{\text{res},1} &= \sqrt{F_{21}^{2} + F_{31}^{2}} \\
            &= \sqrt{(1{,}1238 \, \mathrm{N})^{2} + (1{,}4983 \, \mathrm{N})^{2}} \\
            &= \sqrt{1{,}2628\, \mathrm{N}^2 + 2{,}2450\, \mathrm{N}^2} \\
            &= \sqrt{3{,}5078\, \mathrm{N}^2} \\
            &= 1{,}8729 \, \mathrm{N} \simeq 1{,}9 \, \mathrm{N}
        \end{align*}
        De resulterende kracht is ongeveer $1{,}9 \, \mathrm{N}$ groot.
    \end{solution}
\end{exercise}
%———————————————————————————————————————

%———————————————————————————————————————
% OEFENING EXTRA 2: Vierkant (Complexe Vectoriële Optelling)
\begin{exercise}
    Vier puntladingen bevinden zich op de hoekpunten van een vierkant met zijde $0{,}20 \, \mathrm{m}$.  
    Lading $Q_{1} = +2{,}0 \cdot 10^{-6} \, \mathrm{C}$ bevindt zich linksonder. Lading $Q_{2} = -2{,}0 \cdot 10^{-6} \, \mathrm{C}$ bevindt zich rechtsonder. Lading $Q_{3} = +2{,}0 \cdot 10^{-6} \, \mathrm{C}$ bevindt zich rechtsboven. Lading $Q_{4} = -2{,}0 \cdot 10^{-6} \, \mathrm{C}$ bevindt zich linksboven.
    
    \begin{question}
        Hoe liggen deze drie krachten ten opzichte van elkaar?
        \begin{hint}
            Maak eerst een snelle tekening van de drie afzonderlijke krachten die op $Q_{1}$ inwerken ($\vec{\mathbf{F}}_{21}$, $\vec{\mathbf{F}}_{31}$ en $\vec{\mathbf{F}}_{41}$).    
        \end{hint}
        \begin{multipleChoice}
            \choice{Ze wijzen alle drie naar rechtsboven.}
            \choice[correct]{Twee krachten staan loodrecht op elkaar, de derde ligt exact op dezelfde as (diagonaal) als hun resulterende kracht.}
            \choice{Ze heffen elkaar perfect op tot een netto kracht van 0 N.}
        \end{multipleChoice}
        
        \begin{hint}
            $Q_2$ en $Q_4$ zijn negatief, dus zij trekken het positieve $Q_1$ aan. Dit vormt twee loodrechte krachten (langs de zijden van het vierkant). $Q_3$ is óók positief, dus die stoot $Q_1$ af. Die kracht ligt op de verlengde diagonaal van het vierkant!
        \end{hint}
    \end{question}
    
    \begin{question}
        Bereken nu de grootte van de totale resulterende kracht op $Q_{1}$.
        
        \[ F_{\text{res},1} = \answer[tolerance=0.05,onlineshowanswerbutton]{0.82} \, \mathrm{N} \]

        \begin{hint}
            Stap 1: Bereken $F_{21}$ en $F_{41}$ met de wet van Coulomb. Omdat ze loodrecht op elkaar staan, kan je hiervan eenvoudig een eerste resultante $\vec{\mathbf{F}}_{24}$ berekenen met Pythagoras.
        \end{hint}
        \begin{hint}
            Stap 2: Bereken $F_{31}$. Let op: de afstand tot $Q_3$ is de diagonaal van het vierkant. Deze diagonaal bereken je met Pythagoras: $r_{13} = \sqrt{(0{,}20\,\mathrm{m})^2 + (0{,}20\,\mathrm{m})^2}$.
        \end{hint}
        \begin{hint}
            Stap 3: $\vec{\mathbf{F}}_{24}$ trekt $Q_1$ naar het midden van het vierkant. $\vec{\mathbf{F}}_{31}$ duwt $Q_1$ weg van het midden (afstoting). Omdat deze vectoren op perfect dezelfde as (de diagonaal) liggen, mag je ze wiskundig gewoon van elkaar aftrekken!
        \end{hint}

        \begin{solution}\nl
            \underline{Gegevens:}
            \begin{align*}
                & \lvert Q_{1} \rvert = \lvert Q_{2} \rvert = \lvert Q_{3} \rvert = \lvert Q_{4} \rvert = 2{,}0 \cdot 10^{-6} \, \mathrm{C} \\
                & r_{12} = r_{14} = 0{,}20 \, \mathrm{m}
            \end{align*}
    
            \underline{Gevraagd:} $F_{\text{res},1}$\\ 
    
            \underline{Oplossing:}\\
            Er werken drie krachten op $Q_1$. We rekenen ze in stappen uit.
            
            \textbf{Stap 1: De krachten langs de randen}\\
            De ladingen $Q_2$ en $Q_4$ zijn negatief. Zij trekken het positieve $Q_1$ aan met krachten $\vec{\mathbf{F}}_{21}$ en $\vec{\mathbf{F}}_{41}$. Omdat afstanden en ladingen identiek zijn, zijn deze krachten even groot:
            \[
                F_{21} = F_{41} = \frac{8{,}99  \cdot 10^9 \, \mathrm{\frac{N \cdot m^2}{C^2}}\cdot (2{,}0 \cdot 10^{-6})\,\mathrm{C} \cdot (2{,}0 \cdot 10^{-6}\,\mathrm{C})}{(0{,}20\,\mathrm{m})^2} = 0{,}899 \, \mathrm{N}
            \]
            We bundelen deze twee loodrechte vectoren met de stelling van Pythagoras tot één resulterende 'trekkracht' ($\vec{\mathbf{F}}_{24}$) in de richting van het midden van het vierkant:
            \[
                F_{24} = \sqrt{(0{,}899\,\mathrm{N})^2 + (0{,}899\,\mathrm{N})^2} = \sqrt{2} \cdot 0{,}899\,\mathrm{N} \simeq 1{,}271 \, \mathrm{N}
            \]
            
            \textbf{Stap 2: De afstotende diagonale kracht}\\
            Lading $Q_3$ bevindt zich op de diagonaal. We berekenen eerst de afstand $r_{13}$ (de lengte van de diagonaal):
            \[
                r_{13} = \sqrt{(0{,}20\,\mathrm{m})^2 + (0{,}20\,\mathrm{m})^2} = \sqrt{0{,}08} \simeq 0{,}2828 \, \mathrm{m}
            \]
            Beide ladingen zijn positief, dus $Q_3$ stoot $Q_1$ af, weg van het centrum van het vierkant:
            \[
                F_{31} = \frac{8{,}99 \cdot 10^9\, \mathrm{\frac{N \cdot m^2}{C^2}} \cdot (2{,}0 \cdot 10^{-6}\,\mathrm{C}) \cdot (2{,}0 \cdot 10^{-6}\,\mathrm{C})}{(\sqrt{0{,}08}\,\mathrm{m})^2} = \frac{0{,}03596}{0{,}08}\,\mathrm{N} = 0{,}4495 \, \mathrm{N}
            \]

            \textbf{Stap 3: Alles samenbrengen}\\
            De aantrekkingskracht $F_{24}$ wijst via de diagonaal \emph{naar het midden}.
            De afstotingskracht $F_{31}$ wijst langs exact dezelfde as \emph{van het midden weg}.
            Omdat ze een tegengestelde zin hebben op dezelfde lijn, mogen we ze simpelweg aftrekken:
            \[
                F_{\text{res}} = F_{24} - F_{31} = 1{,}271 \, \mathrm{N} - 0{,}4495 \, \mathrm{N} = 0{,}8215 \, \mathrm{N} \simeq 8{,}2 \cdot 10^{-1} \, \mathrm{N}
            \]
            De netto kracht trekt lading $Q_1$ schuin naar het centrum van het vierkant.
        \end{solution}
    \end{question}
\end{exercise}
%———————————————————————————————————————

\end{document}